\chapter{Methodology}

This chapter outlines the methodology used in this thesis. First XXXXX. 

\section{Problem Formulation}

The cooperative striker–jammer engagement is modelled as a
Dec-POMDP $\langle \mathcal{N}, \mathcal{S}, \{\mathcal{A}_i\},
\{\mathcal{O}_i\}, P, \mathcal{R}, \gamma\rangle$ in a two-dimensional
plane. The agent set $\mathcal{N} = \mathcal{N}_s \cup \mathcal{N}_j$
is partitioned into $N_s$ \emph{strikers} and $N_j$ \emph{jammers},
operating in an environment containing a set of radars $\mathcal{R}$
and a set of protected targets $\mathcal{T}$. The team objective is to
destroy all targets in $\mathcal{T}$; jammers degrade adversary radar
coverage to clear penetration corridors for the strikers. 

\subsection{Agents}

Each agent $i$ has continuous state $s_i = [x_i, y_i, v_i, \psi_i,
\dot{\psi}_i]^{\top}$ and follows unicycle kinematics with
piecewise-constant acceleration commands.

\subsubsection{Observation Space}
Agents are partially observed. Each agent has a sensing radius
$R_{\text{obs}}$ inside which it directly detects other agents,
radars, and targets. The local observation
$o_i \in \mathcal{O}_i$ is decomposed into four entity-typed channels:
a fixed-size ego vector
$o_i^{\text{self}} = (x_i, y_i, v_i, \psi_i, \dot{\psi}_i, \tau)$
with $\tau$ the normalised episode time, and three variable-length
sets containing relative bearing, range, and entity-specific features
for visible teammates, radars, and targets, respectively.\footnote{
Per-entity feature dimensions are
$d_{\text{agent}}=4$, $d_{\text{radar}}=3$, $d_{\text{target}}=3$;
masked entries are zeroed.}

\subsubsection{Communication}
Agents share observations over a range-limited, multi-hop graph
following \cite{Wang2025CooperativeUAV}. The set of agents in
\emph{direct} communication with $i$ is
\begin{equation}
\mathcal{N}_i^{1} = \bigl\{ i' \in \mathcal{N} \setminus \{i\}
\;\bigl|\; \rho_{ii'} \le R_{\text{comm}} \bigr\},
\label{eq:comm_direct}
\end{equation}
where $\rho_{ii'}$ is the inter-agent distance. The full
communication subgroup $\mathcal{S}_i$ is obtained by transitive
closure (multi-hop relaying):
\begin{equation}
\mathcal{S}_i = \{i\} \cup \mathcal{N}_i, \qquad
\mathcal{N}_i = \bigl\{ i'' \;\bigl|\; \exists\,
i'' \xleftrightarrow{\mathcal{N}^1} i \bigr\}.
\label{eq:comm_multihop}
\end{equation}
Members of $\mathcal{S}_i$ pool their local detections, so the set of
entities visible to $i$ for entity type
$\xi \in \{\text{agent}, \text{radar}, \text{target}\}$ is
\begin{equation}
\mathcal{D}^{\xi}_{\mathcal{S}_i} = \bigcup_{i' \in \mathcal{S}_i}
\mathcal{D}^{\xi}_{i'},
\qquad
\mathcal{D}^{\xi}_{i'} = \bigl\{ e \in \xi \;\bigl|\;
\rho_{i'e} \le R_{\text{obs}} \bigr\}.
\label{eq:shared_obs}
\end{equation}
Known prior entities are always included; only unknown entities
require detection through \eqref{eq:shared_obs}.

\subsubsection{Action Space}
Each agent $i$ commands a tangential acceleration and a yaw-rate
acceleration through two independent categorical heads, yielding the
factored action $a_i = (a_i^{v}, a_i^{\omega})$ with
$a_i^{v}, a_i^{\omega} \in \{0, 1, \dots, K{-}1\}$, $K=7$. Each
discrete index maps to a multiplier
$c_k \in \{-1,\,-0.5,\,-0.1,\,0,\,+0.1,\,+0.5,\,+1\}$ scaling the
maximum allowable change:
\begin{equation}
v_{i,t+1} = \mathrm{clip}\!\left(v_{i,t} +
c_{a_i^{v}} \cdot \Delta v_{\max},\;v_{\min},\,v_{\max}\right),
\label{eq:vel_update}
\end{equation}
\begin{equation}
\dot{\psi}_{i,t+1} = \mathrm{clip}\!\left(\dot{\psi}_{i,t} +
c_{a_i^{\omega}} \cdot \Delta \dot{\psi}_{\max},\;
-\dot{\psi}_{\max},\,\dot{\psi}_{\max}\right).
\label{eq:omega_update}
\end{equation}
The joint policy is therefore a product of two role-specific
multi-categorical distributions, one shared across all strikers and
one shared across all jammers (parameter sharing within roles).

\paragraph{Strike effect.}
Each striker $i \in \mathcal{N}_s$ has a forward engagement region
defined by a range $r^{\text{Stk}}$ and a half-angle
$\tfrac{1}{2}\alpha^{\text{Stk}}$ centred on its heading $\psi_i$.
A target $j \in \mathcal{T}$ is engaged whenever
\begin{equation}
\mathcal{A}_i^{\mathcal{T}} = \Bigl\{ j \;\Bigl|\;
\rho_{ji} \le r^{\text{Stk}} \;\wedge\;
|\theta_{ji}| \le \tfrac{1}{2}\alpha^{\text{Stk}} \Bigr\},
\label{eq:strike_set}
\end{equation}
where $\theta_{ji}$ is the bearing of $j$ in the body frame of $i$.
Engagement is deterministic: any $j \in \mathcal{A}_i^{\mathcal{T}}$
is removed from $\mathcal{T}$.

\paragraph{Jamming effect.}
Each alive jammer $i \in \mathcal{N}_j$ is assumed to continuously
suppress every adversary radar irrespective of range. The effective
detection range of radar $r \in \mathcal{R}$ is therefore
\begin{equation}
R^{\text{eff}}_r =
\begin{cases}
\max\!\bigl(0,\, R_r - \Delta R\bigr) & \text{if } |\mathcal{N}_j^{\text{alive}}| \ge 1, \\[2pt]
R_r & \text{otherwise},
\end{cases}
\label{eq:jam_effect}
\end{equation}
with $R_r$ the nominal radar range and $\Delta R$ the jamming-induced
range reduction.

\subsection{Environment}


\subsubsection{Low Fidelity Model}


\subsubsection{Higher Fidelity Model}
The radar detection performance is modeled using the standard monostatic radar equation combined with a simplified burn-through condition for jamming. Three effective detection ranges are defined: the unconstrained range, the main-lobe jammed range, and the side-lobe jammed range.

The unconstrained radar range is obtained by equating the received target echo power to the receiver noise power:
\begin{equation}
R_{\text{unconstrained}} =
\left(
\frac{P_t G_t G_r \lambda^2 \sigma}
{(4\pi)^3 k T_0 B_n L}
\right)^{1/4}
\end{equation}

Jamming is modeled using a burn-through condition where the received jammer power equals the target echo power. For main-lobe jamming, the effective detection range is
\begin{equation}
R_{\text{main}} =
\left(
\frac{P_t G_t \sigma}{4\pi P_J G_J} R_J^2
\right)^{1/4}
\end{equation}

For side-lobe jamming, accounting for reduced antenna gain toward the jammer, the effective range becomes
\begin{equation}
R_{\text{side}} =
\left(
\frac{\sigma}{4\pi}
\cdot
\frac{P_t G_t}{P_J G_J}
\cdot
\frac{G_M R_J^2}{G_S}
\right)^{1/4}
\end{equation}

For a monostatic radar, the main-lobe receive gain satisfies $G_M = G_t$.

In these expressions, $P_t$ and $P_J$ denote the radar and jammer transmit powers, $G_t$, $G_r$, and $G_J$ are the corresponding antenna gains, $\lambda$ is the wavelength, $\sigma$ is the target radar cross section, $R_J$ is the radar-to-jammer distance, $G_S$ is the radar side-lobe gain, $k$ is Boltzmann's constant, $T_0$ is the system temperature, $B_n$ is the receiver bandwidth, and $L$ represents system losses.

The radar coverage is then defined as a piecewise function of angle, where the detection range is reduced to $R_{\text{main}}$ or $R_{\text{side}}$ within the jammer sector, and remains $R_{\text{unconstrained}}$ elsewhere.


\section{Methods}


\subsection{MAPPO}
We optimise the clipped PPO objective~\cite{schulman2017ppo} in the MAPPO
variant~\cite{yu2021mappo}, with parameter sharing inside each role:
$\pi_{\theta_\rho}$ for role $\rho\in\{s,j\}$, and role-specific centralised
critics $V_{\phi_\rho}$.

\paragraph{Advantage estimation.}
With rollouts of length $T$ collected from $B$ parallel environments, per-role
advantages are computed by GAE~\cite{schulman2016gae}:
\begin{equation}
\delta_t^{\rho} = r_t^{\rho} + \gamma\, V_{\phi_\rho}(s_{t+1})(1-d_t) - V_{\phi_\rho}(s_t),
\qquad
\hat{A}_t^{\rho} = \sum_{l=0}^{T-t-1}(\gamma\lambda)^{l}\,\delta_{t+l}^{\rho},
\end{equation}
with returns $\hat{R}_t^\rho=\hat{A}_t^\rho+V_{\phi_\rho}(s_t)$.
We use $\gamma=0.99$, $\lambda=0.95$.

\paragraph{Clipped policy loss.}
For each role, letting
$r_t(\theta)=\exp\!\big(\log\pi_{\theta_\rho}(\bm{a}_t\mid o_t)-\log\pi_{\theta_{\rho,\text{old}}}(\bm{a}_t\mid o_t)\big)$,
\begin{equation}
\mathcal{L}^{\text{PPO}}_\rho(\theta)
= -\,\mathbb{E}_t\!\left[\min\!\big(r_t(\theta)\tilde{A}_t^\rho,\;
    \mathrm{clip}(r_t(\theta),1\!-\!\epsilon,1\!+\!\epsilon)\,\tilde{A}_t^\rho\big)\right]
- \beta\,\mathbb{E}_t[\mathcal{H}(\pi_{\theta_\rho})],
\end{equation}
with clip range $\epsilon=0.2$ and entropy coefficient $\beta=0.01$.
$\tilde{A}_t^\rho$ denotes the per-minibatch standardised advantage,
$\tilde{A}=(\hat{A}-\mu_{\hat A})/\sigma_{\hat A}$.

\paragraph{Value loss.}
\begin{equation}
\mathcal{L}^{V}_\rho(\phi)=\mathbb{E}_t\!\left[\big(V_{\phi_\rho}(s_t)-\hat{R}_t^\rho\big)^{2}\right].
\end{equation}

\paragraph{Reward normalisation.}
A running second moment over per-step team-mean rewards $\bar r_t=\tfrac{1}{|\mathcal N|}\sum_i r_{i,t}$
is tracked with a Welford estimator; rewards are then divided (no mean
subtraction) by $\max(\sigma,\varepsilon)$ with $\varepsilon=10^{-2}$
before advantage computation.

\paragraph{Adaptations from standard MAPPO.}
\begin{enumerate}[leftmargin=1.25em]
\item \emph{Role-factorised critic.} Instead of a single shared centralised
critic, we instantiate one actor $\pi_{\theta_\rho}$ \emph{and} one critic
$V_{\phi_\rho}$ per role, each optimised by a dedicated Adam instance. The
critics share the same global-state input but learn role-specific value
functions, which we found necessary because the two roles receive different
reward mixtures (Sec.~\ref{subsec:rewards}).
\item \emph{Unclipped value loss.} We omit the PPO value clip; the plain MSE
in~(4) was empirically stable in our setup and avoided a known pathology
where clip-range mis-calibration freezes the critic early in training.
\item \emph{Std-only reward scaling.} We divide rewards by the running std
without mean subtraction, preserving the sign/zero of sparse terminal
rewards (kill / mission-complete), which would otherwise drift to a
non-zero baseline under a Welford-mean subtraction.
\item \emph{Factorised MultiCategorical} action distribution with logit
sanitisation $\mathrm{clip}(\ell,-20,20)$ and NaN guarding, to tolerate
occasional extreme logits produced during the early exploration phase.
\item \emph{Per-role gradient clipping.} Actor/critic gradients are clipped
independently at $\lVert g\rVert_2\le 1$ before each optimiser step.
\end{enumerate}

\subsection{Flexible Observation Feature Encoding (FOFE)}
The observation encoder replaces a fixed top-$K$ flat vector with a
permutation-invariant set encoder applied per observation channel, following
the Sequence-Exchangeable Encoding (SEE) block of Wang et al.~\cite{Wang2025CooperativeUAV}.
Each channel $c\in\{\text{ag},\text{tg},\text{rd}\}$ provides a variable-size
entity set $\mathcal{X}^{(c)}\!=\!\{x_1,\dots,x_n\}$ with a visibility mask
$m\in\{0,1\}^n$.

\paragraph{SEE layer.}
For input $X\in\mathbb{R}^{n\times d_{\text{in}}}$ and hidden width $h$,
\begin{align}
\tilde X &= \mathrm{LayerNorm}(X),\\
E &= \tilde X W_e,\qquad Z=\tilde X W_g,\qquad G=\mathrm{SiLU}(Z),\\
E_{\text{agg}} &= \max_{i:\,m_i=1} E_{i,:} \;\in \mathbb{R}^{h} \quad\text{(masked max-pool),}\\
\mathrm{SEE}(X,m)_{i,:} &= \big[\,G_{i,:}\odot E_{i,:}\,\big\Vert\,G_{i,:}\odot E_{\text{agg}}\,\big] \in \mathbb{R}^{2h},
\end{align}
with layer-norm omitted in the first SEE block. The left branch carries
local per-entity context, the right branch injects a gated global summary;
masked rows are zeroed after the concat.

\paragraph{FOFE block.}
A channel encoder stacks $L$ SEE layers, applies a shared per-entity MLP
$f_\psi$, then collapses the set by masked max-pool:
\begin{equation}
    \mathrm{FOFE}(\mathcal X,m)
    = \max_{i:\,m_i=1} f_\psi\!\big(\mathrm{SEE}_L\circ\cdots\circ\mathrm{SEE}_1(X,m)\big)_i
    \;\in \mathbb{R}^{D_{\mathrm{fofe}}}.
\end{equation}
By construction this is permutation-invariant and agnostic to the number of
visible entities.

\paragraph{Actor network (per role).}
The self-observation $o^{\text{self}}\in \mathbb{R}^{6}$ (ego kinematics + alive
flag) bypasses FOFE and is mapped by an MLP. The three entity channels are
encoded independently and concatenated with the self-embedding before a
fusion MLP and a factorised action head:
\begin{equation}
    z = \big[\,g_\xi(o^{\text{self}})\,\big\Vert\,\mathrm{FOFE}_{\text{ag}}\,\big\Vert\,
            \mathrm{FOFE}_{\text{tg}}\,\big\Vert\,\mathrm{FOFE}_{\text{rd}}\,\big],\qquad
    \bm{\ell} = W_{\text{head}}\,\mathrm{FusionMLP}(z).
\end{equation}

\paragraph{Critic network (per role).}
The critic receives the decomposed \emph{global} state: all agents
($d{=}7$: pos, speed, heading, $\omega$, role, alive), all targets ($d{=}3$),
all radars ($d{=}4$: pos, active flag, effective range), plus a scalar
time feature $\tau\in[0,1]$:
\begin{equation}
    V_{\phi_\rho}(s) = W_v\;\mathrm{FusionMLP}\!\Big(
    \big[\,\mathrm{FOFE}_{\text{ag}}(s)\,\big\Vert\,\mathrm{FOFE}_{\text{tg}}(s)\,\big\Vert\,
        \mathrm{FOFE}_{\text{rd}}(s)\,\big\Vert\,\tau\,\big]\Big).
\end{equation}

\paragraph{Adaptations from standard FOFE.}
\begin{enumerate}[leftmargin=1.25em]
\item \emph{Per-channel decomposition.} Wang's original encoder processes a
single heterogeneous entity stream. We split the observation into
semantically typed channels (agents, targets, radars) each with an
independent FOFE block, and treat the ego state as a separate fixed-size
MLP branch. This enforces a type-aware prior and makes channel-level
diagnostics (dominance, collapse rate) tractable.
\item \emph{Centralised FOFE critic with time feature.} The critic
consumes the full global state as three entity sets \emph{plus} an explicit
normalised-time scalar $\tau$, enabling the baseline to discriminate early
from late episode phases without extending the set with a dummy entity.
\item \emph{Masked max-pool with all-masked fallback.} When no entity in a
channel is visible ($\sum_i m_i=0$) we substitute a zero vector rather
than propagating $-\infty$ from the masked max, ensuring stable gradients
under aggressive occlusion.
\item \emph{Role-specific FOFE weights.} Striker and jammer each own a full
set of three FOFE blocks (actor and critic); no weight sharing across roles.
\end{enumerate}


\subsection{Reward Structure}
\label{subsec:rewards}
Rewards are applied only to alive agents. In the FOFE setup used here
($1$ striker, $1$ jammer, $1$ target, $1$ radar), the implemented reward
contains two sparse event terms, one terminal bonus, and four dense shaping
terms plus a control penalty.

\paragraph{Sparse event rewards.}
For an event $e$, agent $i$ receives
\begin{equation}
r_i^{e} = \eta\,\bar R_e + (1-\eta)\,\tilde R_i^{e},
\end{equation}
where $\bar R_e$ is the event reward shared equally across alive agents and
$\tilde R_i^{e}$ is the individual credited part. A target kill rewards
successful engagement, an agent loss penalises radar exposure, and a
terminal bonus is given when all targets are destroyed.

\paragraph{Dense shaping terms.}
The per-step reward is
\begin{equation}
r_{i,t} = r_i^{e} + r^{\text{time}}_{i,t} + r^{\text{border}}_{i,t}
+ r^{\text{radar}}_{i,t} + r^{\text{strike}}_{i,t}
+ r^{\text{form}}_{i,t} + r^{\text{ctrl}}_{i,t}.
\end{equation}
The active dense terms are
\begin{align}
    r^{\text{time}}_{i,t}   &= -c_{\text{time}}, \\
    r^{\text{border}}_{i,t} &= -\phi_{\text{border}}(d^{\text{edge}}_{i,t}), \\
    r^{\text{radar}}_{i,t}  &= -\sum_{k\in\mathcal R}\phi_{\text{radar}}(d^{\text{radar}}_{ik,t}), \\
    r^{\text{strike}}_{i,t} &= -\phi_{\text{target}}(d^{\text{target}}_{i,t}),
    \qquad i\in\mathcal N_{\text{str}}, \\
    r^{\text{form}}_{i,t}   &= -\phi_{\text{form}}(d^{\text{ally}}_{i,t}),
    \qquad i\in\mathcal N_{\text{jam}}, \\
    r^{\text{ctrl}}_{i,t}   &= -\psi(u_{i,t}).
\end{align}
$r^{\text{time}}$ discourages unnecessarily long trajectories.
$r^{\text{border}}$ penalises flying close to the map boundary.
$r^{\text{radar}}$ penalises entering radar threat regions.
$r^{\text{strike}}$ pulls the striker toward the nearest alive target and is
zero at contact. $r^{\text{form}}$ keeps the jammer close to the striker.
$r^{\text{ctrl}}$ penalises aggressive acceleration and heading changes.

\begin{table}[t]
\centering
\caption{Reward parameters used in the FOFE setup. Inactive code options
(jammer-approach, progress rewards, jam bonus, same-role separation,
striker-formation reward, and paper mission reward) are omitted.}
\label{tab:reward_params_fofe}
\begin{tabular}{@{}lll@{}}
\toprule
Component & Parameter(s) & Value \\
\midrule
Target kill & $R_{\text{kill}}$ & $+3$ \\
Agent destroyed & $R_{\text{loss}}$ & $-3$ \\
Mission complete & $R_{\text{term}}$ & $+1$ \\
Team sharing & $\eta$ & $0.5$ \\
Time penalty & $c_{\text{time}}$ & $0.01$ \\
Border penalty & $w_{\text{border}},\,d^{\max}_{\text{border}}$ & $0.05,\;0.05$ \\
Radar penalty & $w_{\text{radar}},\,d^{\max}_{\text{radar}}$ & $0.05,\;0.10$ \\
Striker target penalty & $w_{\text{target}},\,d^{\max}_{\text{target}}$ & $0.10,\;1.00$ \\
Jammer formation penalty & $w_{\text{form}},\,d^{\text{ref}}_{\text{form}}$ & $0.10,\;0.50$ \\
Control effort & $\lambda_a,\,\lambda_\omega$ & $0.01,\;0.01$ \\
\bottomrule
\end{tabular}
\end{table}



\section{Performance Metrics}

\subsection{KPIs}


\subsubsection{Emergent Tactics}

What are possible tactics for the agent to take? How can these be quantized? 

